

\documentclass[a4paper,12pt]{article}
\usepackage{geometry}
\geometry{a4paper,left=18mm,right=18mm, top=2cm, bottom=2cm}
 

\usepackage{lmodern}
\usepackage[T1]{fontenc}
\usepackage[utf8]{inputenc}
\usepackage[ngerman]{babel}
\usepackage[solution_on, random=0, info_off]{srdp-mathematik} % solution_on/off, random, info_on/off
%
\usepackage{blindtext}
%
%\pagestyle{plain} %PAGESTYLE: empty, plain
%\onehalfspacing %Zeilenabstand
%\setcounter{secnumdepth}{-1} % keine Nummerierung der Überschriften
%
%
%%%%%%%%%%%%%%%%%%%%%%%%%%%%%%%%%%%%%%%%%%%%%%%%%%%%%%%%%%%%%%%%%
%%%%%%%%%%%%%%%%%%%%%% DOKUMENT - ANFANG %%%%%%%%%%%%%%%%%%%%%%%%
%%%%%%%%%%%%%%%%%%%%%%%%%%%%%%%%%%%%%%%%%%%%%%%%%%%%%%%%%%%%%%%%%
%
%

%\NewEnviron{test}{
%	\begin{enumerate}
%		\item {\BODY}
%	\end{enumerate}
%}

\begin{document}

\begin{titlepage}
\begin{flushright}
\includegraphics[width=0.3\textwidth]{logo.eps}

\Huge \textsc{1. Mathematikschularbeit.} \\

\textsc{\Large am Donnerstag, 16. Dezember 2021} \\

\textsc{\Large Klasse 6c} \\ [0.5cm]

\Large Name: \rule{8cm}{0.4pt} \\

\Large \textsc{Viel Erfolg!!}
\end{flushright}

\vfil 

\beurteilungsraster{0.85}{0.68}{0.5}{1/3}{ % Prozentschluessel
				T1={24}, 				% Punkte im Teil 1  
				AP={4}, 				% Ausgleichspunkte aus Teil 2  
				T2={20}, 				% Punkte im Teil 2
				}
%\normalsize
%Die Schularbeit besteht aus Teil 1 (Grundkompetenzen) und Teil 2 (vernetzte und erweiterte Grundkompetenzen). In Teil 1 sind 32 und in Teil 2 sind 16 Punkte zu erreichen.\\ 
%Für eine positive Beurteilung gilt: mindestens 24 Punkte insgesamt!
%
%
%\leer
%
%\renewcommand{\arraystretch}{1.5}
%\begin{center}
%\begin{tabular}{p{5cm}p{4cm}l}
%&Punkte&davon erreicht\\ \hline
%Teil 1&32&\\\hline
%Teil 2  &16 &\\\hline
%& $\Sigma= 48$ & \\\hline
%\end{tabular}
%\vfil
%Beurteilung: \rule{8cm}{0.4pt}
%\end{center}
%\renewcommand{\arraystretch}{1}


\end{titlepage}
				
				
				
dsalkdöask
dlksaöd


\end{document}


% Aufgabenliste: 4, 10